\uuid{pfKZ}
\titre{Espace des polyômes}
\auteur{hanani ; mzari}
\organisation{exo7}
\licence{CC BY-NC-SA 4.0}
\datecreate{2020-01-20}
\niveau{L1}
\module{Algèbre}
\chapitre{Espace vectoriel}
\sousChapitre{Définition, sous-espace}
\difficulte{3}
\theme{Espaces vectoriels, Espace des polyômes}
\exo7qcmid{lille-2.espaces-vectoriels.015}

\contenu{
\begin{question}
\qclassification{106.01}
\qtags{motcle=Espace des polyômes}
Soit $n$ un entier $\ge 1$ et  $E=\{P \in \Rr[X] \, ; \, \deg P< n\}$, muni des opérations usuelles. Quelles sont les assertions vraies ?
\begin{answers}
\bad{$0\notin E$.}
\good{$E$ est stable par addition.}
\good{$E$ est stable par multiplication par un scalaire.}
\bad{$E$ n'est pas un espace vectoriel.}
\end{answers}
\begin{explanations} On a : $0\in E$ et on vérifie que $E$ est stable par addition et multiplication par un scalaire. Donc $E$ est un sous-espace vectoriel de $\Rr[X]$.
\end{explanations}
\end{question}
}
